\documentclass[11pt]{article}
\usepackage[OT1]{fontenc}
\renewcommand{\rmdefault}{cmr}
\renewcommand{\sfdefault}{cmss}
\renewcommand{\ttdefault}{cmtt}
\usepackage{amsmath,amssymb,amsthm,mathtools}
\usepackage[margin=27mm]{geometry}
\usepackage{microtype,booktabs,longtable,array}
\usepackage[hidelinks]{hyperref}
\setlength{\emergencystretch}{2em}
\setlength{\parskip}{3pt plus 1pt}
\setcounter{tocdepth}{1}
\newtheorem{theorem}{Theorem}[section]
\newtheorem{lemma}[theorem]{Lemma}
\newtheorem{proposition}[theorem]{Proposition}
\newtheorem{corollary}[theorem]{Corollary}
\theoremstyle{definition}
\newtheorem{definition}[theorem]{Definition}
\newtheorem{example}[theorem]{Example}
\theoremstyle{remark}
\newtheorem{remark}[theorem]{Remark}
\newcommand{\N}{\mathbb N}
\newcommand{\Q}{\mathbb Q}
\newcommand{\R}{\mathbb R}
\newcommand{\Z}{\mathbb Z}
\newcommand{\supp}{\operatorname{supp}}
\newcommand{\Con}{\operatorname{Con}}
\newcommand{\ord}{\operatorname{ord}}
\newcommand{\dom}{\operatorname{dom}}
\newcommand{\sgn}{\operatorname{sgn}}
\newcommand{\res}{\operatorname{Res}}
\newcommand{\rad}{\operatorname{rad}}
\newcommand{\CK}{\mathsf{CK}}
\newcommand{\KT}{\mathsf{KT}}
\newcommand{\DCB}{\mathsf{DCB}}
\newcommand{\FS}{\mathsf{FS}}
\newcommand{\Ag}{\mathsf{Ag}}
\newcommand{\Know}{\mathsf K}
\newcommand{\Truth}{\mathsf T}
\newcommand{\Prv}{\mathsf{Pr}}
\newcommand{\code}[1]{\ulcorner #1\urcorner}
\newcommand{\Tot}{\operatorname{Tot}}
\numberwithin{equation}{section}

\title{The Domain of Hyperreal Differentiation\ Is Always Proper}
\author{}
\date{28 October 2026\\Working manuscript}
\begin{document}
\maketitle
\begin{abstract}
Alexander and Dawson define differentiation on those elements of an
ultrapower of the reals which have a representative differentiable on an
ultrafilter-large set. For an additive idempotent ultrafilter concentrated
immediately to the right of zero, the derivative is independent of the
chosen differentiable representative. They ask whether its domain can be
the whole ultrapower, and whether the answer depends on the ultrafilter.
We prove in ZFC that the domain is proper for every such ultrafilter.
More strongly, for each ultrafilter there is a bounded strictly positive
function such that every equivalent representative is discontinuous on
an ultrafilter-large set. A dichotomy for the rational coordinates of a
Hamel basis supplies the construction. The resulting excluded element
is a positive infinitesimal, so the differentiation domain is not an
order-convex subfield. The witness may depend on the ultrafilter; no
measurability or common-witness assertion is made.
\end{abstract}
\tableofcontents

\section{A derivative of a hyperreal number}

An ultrapower represents a number by a real-valued function. This makes
it tempting to differentiate the function and regard the resulting class
as the derivative of the number. The difficulty is that many functions
represent the same number. Agreement on a large set need not imply
agreement of derivatives, and an equivalence class may contain both
well-behaved and badly behaved functions.

Alexander and Dawson \cite{AD} identify an algebraic condition on the
ultrafilter which resolves the first difficulty. If the ultrafilter is
additively idempotent and concentrates near zero from the right, any two
equivalent representatives whose derivatives exist on large sets have
equivalent derivatives. Their Corollary~10 therefore defines a derivative
on a distinguished domain $\mathcal D_p$. Immediately afterwards they ask
whether every hyperreal has a representative of the required kind, and
whether this depends on $p$.
Admissible ultrafilters exist in ZFC, as recorded in their Lemma~7;
the universal statement below is therefore nonvacuous.

There are two different questions here. It is easy to take a good
representative and alter it away from a large set so that it becomes
badly behaved. That does not exclude its equivalence class from
$\mathcal D_p$: the original good representative still exists. To prove
properness one must construct a class for which \emph{every}
representative fails. The proof below is designed around that quantifier.

\begin{theorem}[Universal properness]\label{thm:main}
Let $p$ be an additive idempotent ultrafilter on all subsets of $\R$
such that $(0,\epsilon)\in p$ for every $\epsilon>0$. There is a function
$f:\R\to(0,1]$ with the following properties.
\begin{enumerate}
\item For every $\eta>0$, $\{x:f(x)<\eta\}\in p$.
\item For every $g:\R\to\R$ with $\{x:g(x)=f(x)\}\in p$, there is a set
      $S_g\in p$ at every point of which $g$ is discontinuous.
\end{enumerate}
Consequently $[f]_p$ is a positive infinitesimal outside $\mathcal D_p$,
and $\mathcal D_p\ne\R^{\R}/p$.
\end{theorem}

The first assertion places the excluded element very close to zero in
the order of the ultrapower. The second is the essential obstruction:
it survives all changes of representative. The proof uses finite
intersections of large sets only. It never asks an ultrafilter to be
closed under countably many intersections.

\section{Ultrafilters, addition and the differentiation domain}

An ultrafilter $p$ on a set $X$ is a family of subsets closed under
supersets and finite intersections, excluding the empty set, and
containing exactly one of $A$ and $X\setminus A$ for each $A\subseteq X$.
Call its members \emph{$p$-large}. It is principal at $x$ if its members
are exactly the sets containing $x$; otherwise it contains no finite set.

For subsets of an additive group, put $A-x=\{y:x+y\in A\}$. Addition of
ultrafilters on the underlying discrete group is defined by
\[
 A\in p+q \quad\Longleftrightarrow\quad
 \{x:A-x\in q\}\in p.
\]
Thus idempotence $p+p=p$ says that whenever $A$ is large, a large set of
translations has a large intersection pattern with $A$:
\begin{equation}\label{eq:star}
 A\in p\quad\Longrightarrow\quad
 A^*:=A\cap\{x:A-x\in p\}\in p.
\end{equation}
Near-zero concentration implies that $p$ is nonprincipal: no single real
belongs to every interval $(0,\epsilon)$.

Two functions $f,g:\R\to\R$ represent the same element of
$\R^{\R}/p$ exactly when $\{x:f(x)=g(x)\}\in p$. Addition,
multiplication and order are defined pointwise modulo $p$. The usual
ultraproduct argument gives an ordered field extending $\R$ through
constant functions.

Let $D_f$ be the set of points where $f$ has an ordinary finite two-sided
derivative. The class $[f']_p$ exists when $D_f\in p$: extend the partial
derivative arbitrarily outside $D_f$. The extension does not affect its
class. The source domain is
\begin{equation}\label{eq:domain}
 \mathcal D_p=\{[f]_p:D_f\in p\}.
\end{equation}
Membership means that \emph{some} representative has this property.
Merely having differentiability points arbitrarily close to zero is
weaker than $D_f\in p$ and does not replace this definition.

For context we give the argument behind the well-definedness result
of \cite[Theorem~9]{AD}.

\begin{proposition}\label{prop:well-defined}
If $[f]_p=[g]_p$ and $D_f,D_g\in p$, then $[f']_p=[g']_p$.
Moreover, $\mathcal D_p$ is a subfield of the ultrapower and
$\partial[f]_p=[f']_p$ is a derivation from it to the ultrapower.
\end{proposition}
\begin{proof}
Let $A=\{f=g\}\cap D_f\cap D_g$. By \eqref{eq:star}, $A^*\in p$.
Fix $x\in A^*$. For every $\delta>0$, the set
$(A-x)\cap(0,\delta)$ is large and hence nonempty. Along the arbitrarily
small positive increments in this set, the difference quotients for
$f$ and $g$ at $x$ are equal, since both $x$ and $x+h$ lie in $A$.
Both ordinary derivatives exist, so their values agree. They therefore
agree on $A^*$, proving the first assertion.

Constants are in the domain. Intersecting differentiability sets and
applying the ordinary sum and product rules proves closure under sums,
negatives and products and proves the derivation laws. If $[f]_p\ne0$,
then $\{f\ne0\}\in p$. Define $h(x)=1/f(x)$ where $f(x)\ne0$ and
$h(x)=0$ otherwise. At every point of $D_f\cap\{f\ne0\}$, continuity
of $f$ makes this the ordinary reciprocal in a neighborhood. Hence $h$
is differentiable there, $[h]_p=[f]_p^{-1}$, and the domain is a field.
\end{proof}

The proposition explains why the properness question is meaningful.
It supplies no reason that every equivalence class has a suitable
representative. The next lemma isolates what will prevent one.

\section{An obstruction which survives a change of representative}

\begin{lemma}\label{lem:obstruction}
Suppose $f:\R\to(0,1]$ has this property: whenever $T\in p$, $x\in T^*$
and $\delta>0$, there exists
\[
 y\in (T-x)\cap(0,\delta)
 \quad\text{such that}\quad f(x+y)<f(x)/2.
\]
Then every $g$ equivalent to $f$ is discontinuous at every point of a
$p$-large set. In particular $[f]_p\notin\mathcal D_p$.
\end{lemma}
\begin{proof}
Take $T=\{g=f\}\in p$. At each fixed $x\in T^*$, $g(x)=f(x)>0$.
The hypothesis gives arbitrarily small positive $y$ with $x+y\in T$
and $g(x+y)=f(x+y)<g(x)/2$. The fixed positive tolerance $g(x)/2$
witnesses discontinuity at $x$. Since $T^*\in p$, the continuity set of
$g$ is not large: otherwise its intersection with $T^*$ would be empty
and large. Differentiability implies continuity, so its differentiability
set cannot be large either. The argument applies to every representative.
\end{proof}

The use of idempotence is visible here. Ordinary equality modulo $p$
only gives the set $T$. Idempotence provides points $x\in T$ at which
the increments leading back into $T$ remain large. Near-zero
concentration then turns those increments into a continuity test.

\section{Rational coordinates and the two cases}

Choose a Hamel basis $B$ for $\R$ as a vector space over $\Q$. For every
$x\in\R$ there is a unique finite expression
\[
 x=\sum_{b\in B}h_b(x)b,\qquad h_b(x)\in\Q.
\]
Each coordinate $h_b:\R\to\Q$ is additive. Write
\[
 \supp(x)=\{b:h_b(x)\ne0\},\qquad \ell(x)=|\supp(x)|.
\]
The support is finite; $\ell(x)=0$ precisely when $x=0$. Disjoint
supports give $\ell(x+y)=\ell(x)+\ell(y)$.

For a function $h:X\to Y$, its pushforward ultrafilter is
$h_*p=\{A\subseteq Y:h^{-1}(A)\in p\}$.

\begin{lemma}\label{lem:push}
If $h:\R\to\Q$ is additive, $h_*p$ is an additive idempotent
ultrafilter on $\Q$. If it is principal, it is principal at zero.
\end{lemma}
\begin{proof}
Put $q=h_*p$. For $A\subseteq\Q$, additivity gives
\begin{align*}
 A\in q+q
 &\Longleftrightarrow
 \{x:h^{-1}(A-h(x))\in p\}\in p\\
 &\Longleftrightarrow
 \{x:h^{-1}(A)-x\in p\}\in p\\
 &\Longleftrightarrow h^{-1}(A)\in p+p=p
 \Longleftrightarrow A\in q.
\end{align*}
The sum of principal ultrafilters at $r$ and $r$ is principal at $2r$.
Idempotence therefore implies $r=2r$, hence $r=0$.
\end{proof}

It follows that precisely one of the following alternatives holds:
some $h_{b*}p$ is nonprincipal; or every $h_{b*}p$ is principal at zero.
The second alternative says $\ker h_b\in p$ for every $b\in B$. It does
not say the intersection of all these kernels is large. That intersection
is just $\{0\}$ and is not large.

\subsection{A nonprincipal coordinate}

Suppose $h=h_b$ has nonprincipal image $q=h_*p$. Enumerate
$\Q=\{r_0,r_1,\ldots\}$ without repetitions, set
$s(r_n)=2^{-n-1}$, and define $f(x)=s(h(x))$.
The function is strictly positive and bounded by $1/2$. For every
$\eta>0$, its upper level set is the inverse image of a finite subset
of $\Q$. A nonprincipal ultrafilter excludes finite sets, so
$\{f<\eta\}\in p$.

Fix $T\in p$ and $x\in T^*$, and put $a=f(x)>0$. The set
$F=\{r:s(r)\ge a/2\}$ is finite. Since $q$ is nonprincipal,
\[
 E_x=\{y:h(y)\notin F-h(x)\}\in p.
\]
For each $\delta>0$, choose
$y\in E_x\cap(T-x)\cap(0,\delta)$. Additivity gives
$h(x+y)\notin F$, hence $f(x+y)<a/2$. Lemma~\ref{lem:obstruction}
applies.

Freeness of the \emph{image} ultrafilter matters. The inverse image of a
finite subset of $\Q$ may be infinite in $\R$; freeness of $p$ alone
would not justify this step.

\subsection{All coordinates vanish modulo the ultrafilter}

Suppose now $\ker h_b\in p$ for every $b$. For a fixed real $x$, put
\[
 K_x=\bigcap_{b\in\supp(x)}\ker h_b.
\]
This is a finite intersection and belongs to $p$. For $y\in K_x$, the
supports of $x$ and $y$ are disjoint.

\begin{lemma}\label{lem:length}
For every $k\in\N$, the set $S_k=\{x:\ell(x)=k\}$ is not in $p$.
Consequently $\{x:\ell(x)>m\}\in p$ for every $m\in\N$.
\end{lemma}
\begin{proof}
$S_0=\{0\}$ is not large. Suppose $k>0$ and $S_k\in p$. For every
$x\in S_k$, the large set $K_x\cap S_k$ is contained in $S_{2k}-x$.
Thus
$S_k\subseteq\{x:S_{2k}-x\in p\}$, so the set on the right is large.
Idempotence gives $S_{2k}\in p+p=p$, contradicting disjointness of
$S_k$ and $S_{2k}$. This proves the first assertion. A finite union of
sets excluded by an ultrafilter is excluded, yielding the second.
\end{proof}

Define $f(x)=2^{-\ell(x)}$. Lemma~\ref{lem:length} gives
$\{f<\eta\}\in p$ for every $\eta>0$. Given $T\in p$ and $x\in T^*$,
write $a=f(x)$. For every $\delta>0$, the set
\[
 (T-x)\cap(0,\delta)\cap K_x\cap\{y:\ell(y)\ge2\}
\]
is large. A member $y$ satisfies
\[
 f(x+y)=2^{-\ell(x)-\ell(y)}\le a/4<a/2.
\]
Again Lemma~\ref{lem:obstruction} applies. These two cases prove
Theorem~\ref{thm:main}.

\section{Consequences and the scope of the answer}

\begin{corollary}
For every admissible $p$, $\mathcal D_p$ is a proper subfield containing
$\R$ and is not order convex in $\R^{\R}/p$.
\end{corollary}
\begin{proof}
The subfield assertion is Proposition~\ref{prop:well-defined}. The main
theorem gives $0<[f]_p<1$ with $[f]_p\notin\mathcal D_p$, whereas the
endpoints $0,1$ are in the domain. This contradicts order convexity.
\end{proof}

Although $[f]_p$ is infinitesimal, it is not zero. Equality in the
ultrapower means equality on a large set, and $f$ is nowhere zero.
Consequently replacing $f$ by the zero function is unavailable. This
distinction between infinitesimal proximity and equality is essential
to the example.

For comparison, the identity function represents a positive infinitesimal
$\Omega=[x\mapsto x]_p$ which does belong to $\mathcal D_p$, with
$\partial\Omega=1$. Thus infinitesimal size alone neither includes
nor excludes an element from the differentiation domain.

The answer to the source's proper-versus-full question is therefore
independent of which admissible $p$ is chosen: it is always proper.
The actual subfields $\mathcal D_p$ may still vary. Our quantifier order
is
\[
 \forall p\ \exists f\ \forall g\sim_p f,
\]
with the additional hypotheses on $p$ stated in the main theorem.
We do not obtain a single function working for all $p$, nor an
identification of domains for different ultrafilters.

The construction uses a Hamel basis and unrestricted functions. It does
not produce Borel or measurable witnesses. Establishing properness with
regular representatives, or determining which regularity classes always
meet $\mathcal D_p$, would be a different question. The proof is also
specific about its hypotheses: additive idempotence produces the sets
$T^*$ and controls rational coordinate images; concentration near zero
produces arbitrarily small increments. Neither is replaced here by
freeness alone.

\paragraph{Preparation.}
This is an AI-assisted hand-proof manuscript prepared for mathematical
review. All arguments used for the properness theorem are included.
No proof-assistant certification or established historical priority is
claimed.

\begin{thebibliography}{9}
\bibitem{AD} S.~A.~Alexander and B.~Dawson,
\emph{Hyperreal differentiation with an idempotent ultrafilter},
Journal of Logic and Analysis 16 (2024), paper 6.
\href{https://doi.org/10.4115/jla.2024.16.6}{doi:10.4115/jla.2024.16.6}.
Source version: \href{https://arxiv.org/abs/2411.14689v1}{arXiv:2411.14689v1},
22 November 2024; especially Definitions~5--6, Theorem~9, Corollary~10
and the question following it.
\end{thebibliography}
\end{document}
